% !TEX program = xelatex
% 请用 XeLaTeX 编译（TeXworks/TeXstudio 里选 XeLaTeX，或命令行：xelatex simulation_methodology.tex）
% 勿用 pdfLaTeX：ctexart + 中文 + Windows 下用 pdfLaTeX 常极慢或看似「一直不出来」。
\documentclass[UTF8,a4paper]{ctexart}
\usepackage{amsmath,amssymb,booktabs,geometry}
\usepackage{hyperref}
\usepackage{bookmark}
\geometry{margin=2.5cm}
\hypersetup{colorlinks=true,linkcolor=blue}

\title{新能源车队调度仿真：建模思路、规模参数与得分含义}
\author{党涵}
\date{}

\begin{document}
\maketitle

\section{总体思路}

题目要求用\textbf{图结构}表示道路并完成\textbf{路径规划}，在动态到达的任务下进行车队调度。本实现采用以下抽象层次。

\subsection{道路网络（图模型）}
\begin{itemize}
  \item 将城市抽象为 $r\times c$ 的\textbf{网格图}：每个交叉点为一个顶点，四邻连接（上下左右）。
  \item 边权为相邻顶点间的\textbf{欧氏距离}（网格轴相邻边权为 $1$）。
  \item 顶点按行主序编号，$0$ 号节点为\textbf{中央仓库（车场）}。
  \item 使用 \textbf{Dijkstra} 单源最短路；从某顶点出发的距离行按需计算并\textbf{缓存}，路径由父指针回溯得到（大图避免 Floyd 的 $O(|V|^3)$）。
\end{itemize}

\subsection{时间推进与任务生成}
\begin{itemize}
  \item 仿真时间为离散时间步，步长 $\Delta t = 0.5$。
  \item 每个时间步以概率 $p_{\mathrm{spawn}} \cdot \Delta t$ 独立生成一个新任务。
  \item 每个任务包含：产生时刻、目标顶点、货物重量（区间内均匀随机）、截止时间（产生时刻加上区间内均匀随机的裕量）。
\end{itemize}

\subsection{车队与能耗}
\begin{itemize}
  \item 每辆车有\textbf{电量上限} $B_{\max}$、\textbf{载重上限} $L_{\max}$；每完成一票配送，车上对应载重释放。
  \item 行驶距离 $d$ 的能耗近似为 $E(d)=\rho\cdot d$，其中 $\rho$ 为\textbf{单位距离能耗}。
  \item 车速用于将距离换算为行驶时间：$t_{\mathrm{travel}} = d / v$。
\end{itemize}

\subsection[基线调度策略]{基线调度策略（\texttt{fleet\_simulation.py}）：仓库多票装车 $+$ 最近邻配送顺序}
车辆\textbf{仅在仓库}且空闲时派发新的一批任务（代码中 \texttt{node} 为仓库且无在送批次）。模型假定\textbf{所有货物在仓库一次装车}，再沿路连续送票，\textbf{本批全部送完后}再执行回仓（与“每送一票必须先回仓”不同）。具体为：
\begin{enumerate}
  \item \textbf{装车（贪心）}：在待分配、已出现且未过期的任务中，按\textbf{重量从大到小}依次装入，直至再加一票会超过 $L_{\max}$，得到本趟\textbf{批次}（可多票）。同重量时排序顺序由实现细节决定（稳定排序下保持候选表原有次序）。
  \item \textbf{配送顺序（最近邻）}：以仓库为起点，在批次内反复选择\textbf{距当前位置图最短路最短}的下一配送点，直至排完访问序（启发式，非最优 TSP）。
  \item \textbf{电量可行性}：计算\textbf{仓库 $\to$ 各配送点（按上述顺序）$\to$ 回仓库} 的总路程能耗；若超过当前电量，则反复从批次中\textbf{去掉当前最轻的一票}并重新最近邻排序，直至总能耗不超过电量，或批次退化为\textbf{仅一票}（仍不足则允许经充电站补能，逻辑见下节）。
  \item \textbf{在途}：按访问序分段行驶；每一段均为“当前位置 $\to$ 下一配送点”，段内仍可采用途经充电站路径。每送完一票仅卸下该票重量；若尚有余票，则\textbf{从当前配送点直接}驶向下一点，无需回仓。
\end{enumerate}
代码中另保留单票选取函数 \texttt{pick\_task\_max\_weight}，便于与“只装一票”等策略做对比，主流程为 \texttt{pick\_batch\_greedy\_max\_weight}。

\subsection[元启发式对比实现]{元启发式对比实现（\texttt{fleet\_metaheuristic.py}）}
在\textbf{与基线相同的} \texttt{SimConfig}、任务生成随机种子及 \texttt{self.\_rng} 消费顺序下，通过子类 \texttt{MetaHeuristicFleetSimulator} 仅替换\textbf{装批平局规则}与\textbf{批次内访问序}；仿真主循环、充电与路径逻辑仍继承基类。元启发式内部随机数使用独立种子（\texttt{seed}+常数偏移），避免干扰任务流。

\paragraph{装车（批次层面）}
仍按\textbf{重量从大到小}再在载重内贪心装入，与基线结构一致；区别为\textbf{同重量}任务按\textbf{截止时间先后}（EDD），再按任务编号 \texttt{tid} 打破平局，对应 \texttt{pick\_batch\_weight\_then\_edd}。这样在不推翻“重货优先装满车”的前提下引入紧迫度信息。

\paragraph{批次内路径优化（元启发式核心）}
派发批次时记录当前仿真时刻 $t_0$，对批次内访问序 $\pi$ 定义\textbf{与总分系数对齐的代理代价}（用于优化，非仿真直接加分公式）：
\begin{equation}
  J(\pi) = \alpha_{d}\, D(\pi) + \alpha_{\mathrm{late}}\, L(\pi).
\end{equation}
其中 $D(\pi)$ 为\textbf{仓库 $\to$ 按 $\pi$ 依次送达各点 $\to$ 回仓库}的图最短路总里程（与电量检验所用回路长度一致）；$L(\pi)$ 为按 $\pi$ 顺序、以 $t_0$ 为起点用行驶时间累加到达各点时，\textbf{正迟到时间}之和，即对每个任务取 $\max\{0,\,t_{\mathrm{arr}}-t_{\mathrm{dl}}\}$ 再求和（$t_{\mathrm{arr}}$ 由连续最短路距离除以车速 $v$ 递推）。

\paragraph{算法流程}
\begin{itemize}
  \item 若批内任务数 $n\le 6$，对全排列枚举，取 $J(\pi)$ 最小的访问序（小规模精确最优）。
  \item 若 $n>6$，采用\textbf{模拟退火}：初解在\textbf{基类最近邻序}与\textbf{按截止时间排序的全序列（EDD）}中取 $J$ 较小者；邻域包含\textbf{交换两位置}、\textbf{反转一段（2-opt 型）}、\textbf{单点移位（Or-opt-1）}；对交换移动辅以\textbf{短禁忌表}，长期无接受则\textbf{升温重启}降温轨迹。
  \item 电量不可行时的\textbf{丢票重排}逻辑与基线相同：仍按最轻票剔除后重新优化访问序（调用同一套批次内优化子程序）。
\end{itemize}

\paragraph{对比实验脚本}
\texttt{fleet\_metaheuristic.py} 中 \texttt{run\_controlled\_comparison()} 逐场景先跑基线 \texttt{FleetSimulator} 再跑子类，并校验任务流签名（产生时刻、目标节点、重量序列）一致，最后输出两套 \texttt{summarize} 得分差，便于报告中的控制变量说明。

\subsection{充电与排队}
\begin{itemize}
  \item 对\textbf{每一段}行驶（含首段出仓、批次内相邻配送点之间、以及送完本批后的回仓段），若沿最短路直达能耗超过当前电量，则先经充电站补能后再继续该段。
  \item 车辆在有载配送过程中也会检查\textbf{动态低电阈值}：阈值随“回仓距离耗电”与“回仓路径拥堵度”上升而提高。电量低于该阈值时，允许先执行一次“当前位置 $\to$ 充电站 $\to$ 原位置”的补能绕行，再继续后续任务。
  \item 每个充电站有有限个\textbf{并发槽位}（代码中为 $2$）。充电用时间区间 $[t_{\mathrm{start}},t_{\mathrm{end}})$ 预约；槽位已满时推进到\textbf{最早能开始充电的时刻}，体现\textbf{排队与负荷}。
\end{itemize}

\section{各规模场景的具体数量}

四组场景除下表所列参数外，还共享：
\[
B_{\max}=500,\quad L_{\max}=250,\quad \rho=0.8,\quad v=1.0,\quad P_{\mathrm{charge}}=75.0,
\]
\[
\alpha_{\mathrm{early}}=10.0,\quad \alpha_{\mathrm{late}}=15.0,\quad \alpha_{d}=0.01.
\]
充电站每站点并发槽位数为 $2$。

\begin{table}[htbp]
  \centering
  \caption[四种仿真规模参数]{四种仿真规模的显式参数（与 \texttt{preset\_scenarios()} 一致）}
  \small
  \begin{tabular}{@{}lcccccccccc@{}}
    \toprule
    名称 & $r$ & $c$ & $|V|=rc$ & 车辆数 & 充电站数 & $T_{\mathrm{sim}}$ &
    $p_{\mathrm{spawn}}$ & 重量区间 & 截止裕量区间 & \texttt{seed} \\
    \midrule
    SMALL      & 25 & 25 & 625   & 10 & 10 & 1000 & 0.40 & $[5,40]$   & $[125,275]$ & 1 \\
    MEDIUM     & 50 & 50 & 2500  & 20 & 20 & 1500 & 0.60 & $[5,75]$   & $[100,300]$ & 2 \\
    LARGE      & 80 & 80 & 6400  & 40 & 40 & 2000 & 0.75 & $[10,125]$ & $[90,250]$ & 3 \\
    XL\_STRESS & 100 & 100 & 10000 & 30 & 25 & 1750 & 1.10 & $[15,150]$ & $[60,175]$ & 4 \\
    \bottomrule
  \end{tabular}
\end{table}

\noindent\textbf{说明：} $|V|$ 为顶点数；边数约为 $2rc-r-c$。
\textbf{XL\_STRESS} 为第四档场景名称：\textbf{XL} 表示网格规模最大（$100\times 100$）；\textbf{STRESS} 表示相对同规模地图\textbf{加压}——任务到达率 $p_{\mathrm{spawn}}$ 更高、截止裕量区间更紧、车辆与充电站数量相对偏少，用于观察调度在拥堵、易超时情形下的表现，与前三档“常规规模”形成对照。

\section{最终得分的含义与计算公式}
\label{sec:score}

最终得分 $S$ 为\textbf{标量评价指标}，对应大作业中“完成越早、路径越短越高分，超时扣分”的定性要求。实现为\textbf{可加}形式。

\subsection{已完成任务的贡献}
设任务在 $t_{\mathrm{fin}}$ 完成，截止 $t_{\mathrm{dl}}$，重量 $w$，本任务相关行驶距离为 $d_{\mathrm{task}}$。
多票同车时，$d_{\mathrm{task}}$ 取\textbf{送达该票的这一小段路}的最短路长度（上一停靠点 $\to$ 该任务顶点，若经充电站则含充电前路段的计程方式与代码 \texttt{travel\_distance} 一致），而非整趟多点的总里程。
\begin{equation}
  \Delta S_{\mathrm{done}} =
  \begin{cases}
    \alpha_{\mathrm{early}}\, w - \alpha_{d}\, d_{\mathrm{task}}, & t_{\mathrm{fin}} \le t_{\mathrm{dl}},\\[0.5em]
    -\alpha_{\mathrm{late}}\,(t_{\mathrm{fin}}-t_{\mathrm{dl}}) - \alpha_{d}\, d_{\mathrm{task}}, & t_{\mathrm{fin}} > t_{\mathrm{dl}}.
  \end{cases}
\end{equation}

\subsection{未完成（过期）任务的惩罚}
待分配任务在仿真中首次被发现已超过截止时间时记一次惩罚：
\begin{equation}
  \Delta S_{\mathrm{expired}} = -\alpha_{\mathrm{late}}\,(t-t_{\mathrm{dl}}).
\end{equation}

\subsection[如何解读得分]{如何解读 $S$}
\begin{itemize}
  \item $S$ \textbf{无绝对物理单位}，是人为设定的效用代理，用于同一套参数下比较策略或随机种子。
  \item \textbf{越大越好}；大量过期或迟到时 $S$ 可显著为负（罚分主导）。
  \item 若修改 $\alpha_{\mathrm{early}},\alpha_{\mathrm{late}},\alpha_{d}$，数值尺度会变，比较时需对齐系数。
\end{itemize}

\section[静态 MILP 与离散仿真对照]{静态 MILP（\texttt{fleet\_gurobi\_static.py}）与离散仿真（\texttt{fleet\_simulation.py}）的对照}
\label{sec:gurobi-compare}

除\textbf{提前已知}全部任务的出现时刻、目标顶点、重量与截止时间（上帝视角）外，仍希望尽可能共用同一“世界”。本节说明\textbf{哪些条件与仿真一致}、\textbf{哪些必然不一致}，避免将 Gurobi 输出与 \texttt{run()} 轨迹做\textbf{逐事件等价}的误读。实现以 \texttt{gurobipy} 求解混合整数规划；具体变量与约束以代码为准。

\subsection{数据与随机世界：可与仿真对齐的部分}

\begin{itemize}
  \item \textbf{场景参数}：与 \texttt{preset\_scenarios()} 中 \texttt{SimConfig} 一致（网格规模、车辆数、充电站数、$T_{\mathrm{sim}}$、$p_{\mathrm{spawn}}$、重量与截止裕量区间、$B_{\max},L_{\max},\rho,v,P_{\mathrm{charge}}$ 及得分系数等）。
  \item \textbf{任务列表与充电站坐标}：代码中通过构造 \texttt{FleetSimulator(cfg)} 后\textbf{仅执行与 \texttt{run()} 相同的任务生成循环}（时间步长 $\Delta t=0.5$，每步以概率 $p_{\mathrm{spawn}}\cdot\Delta t$ 生成一单），\textbf{不执行}车辆派发与 \texttt{step()} 中其余逻辑。由此得到的任务序列与 \texttt{FleetSimulator.\_\_init\_\_} 中已确定的\textbf{仓库顶点}、各充电站\textbf{图节点编号}（\texttt{ChargingStation.node}）与完整仿真在“仅生成任务”意义下一致，从而保证\textbf{最短路距离矩阵}所依赖的端点与仿真相同。
  \item \textbf{图与最短路}：与仿真一样在网格邻接图上用 Dijkstra 得到任意两点间最短路长度；行驶时间仍取 $d/v$。
  \item \textbf{能耗与补电时长}：弧上耗电仍用 $\rho\cdot d$；在充电站将电池从到达电量充至 $B_{\max}$ 所需时间，与仿真中按功率 $P_{\mathrm{charge}}$ 由能量缺口换算的思路一致（代码中为线性约束 \texttt{dur}$\cdot P_{\mathrm{charge}}\ge B_{\max}-B_{\mathrm{arr}}$ 等）。
  \item \textbf{载重}：静态模型中每辆车\textbf{单趟路径}上所承载任务重量之和 $\le L_{\max}$，对应仿真里“在仓库一次装齐本批再出发、本批总重不超过载重上限”的设定。
  \item \textbf{并发槽位数}：每个物理充电站在任意时刻“同时在充”的车辆数上限为 $2$，与仿真里 \texttt{slots}=2 的\textbf{数量}一致（实现机制不同，见下节）。
\end{itemize}

\subsection{模型结构：与仿真不一致或仅为近似的部分}

以下差异来自\textbf{静态优化}与\textbf{离散事件仿真}的不同范式，一般不能仅靠调参消除。

\paragraph{（1）动态多趟 vs.\ 静态单趟}
仿真中车辆在仿真期内可多次回仓、再装新批、再出发，形成\textbf{多条}“出仓—送批—回仓”轨迹；当前 MILP 为\textbf{每辆车一条}从起点仓到终点仓的闭合路径抽象（含可选充电节点），相当于对\textbf{一整段连续作业}做一次性规划，\textbf{不重现}“回仓后再派下一批”的多轮决策过程。因此即便任务全集相同，\textbf{时间轴上的决策次数与仿真也不相同}。

\paragraph{（2）充电排队：连续时间预约 vs.\ 离散时间槽容量}
仿真中充电通过时间区间 $[t_{\mathrm{start}},t_{\mathrm{end}})$ 预约，在槽位占满时按\textbf{最早可开始时刻}推进，是会话列表驱动的\textbf{连续时间}排队。
MILP 中将时间轴按步长 $\Delta_q$（代码参数 \texttt{queue-delta}）分段，用二进制变量刻画充电区间与每个时间槽是否相交，并对每个物理站、每个槽施加\textbf{同时在槽内占用数 $\le 2$}。这是对“两槽并行”的\textbf{离散近似}：\textbf{不能保证}与仿真中每一次 \texttt{\_reserve\_charge}、\texttt{\_next\_charge\_start} 得到的开始/结束时刻逐次相同；$\Delta_q$ 越小通常越细、变量越多，仍未必与连续预约\textbf{事件级等价}。

\paragraph{（3）路径段与失败回退}
仿真中 \texttt{\_begin\_leg\_from\_to} 在电量、充电站可达性等条件下可能失败并触发\textbf{批次回滚}等逻辑；MILP 通常假设所选弧在能量与时间上可行，\textbf{不包含}与仿真完全同构的回滚分支。

\paragraph{（4）目标函数与仿真得分}
仿真得分 $S$ 按第~\ref{sec:score}~节在任务完成、过期等\textbf{离散事件}上累加；MILP 目标为\textbf{线性代理}（如服务任务的重量的按时奖励系数之和、减去迟到项、距离惩罚及未服务惩罚等），用于在静态假设下求较优路径，\textbf{数值上不等于}同一次仿真运行结束后的 $S$，也不宜在无控制实验设计下直接宣称“同分”。

\paragraph{（5）任务子集与规模}
为控制整数规划规模，实现中常只对按重量排序后的\textbf{前若干条}任务建变量（参数 \texttt{max-tasks}）；仿真则对\textbf{全部}到达任务执行在线规则。对比报告时需说明 MILP 求解的是\textbf{子问题}还是经筛选后的实例。

\subsection{对照小结表（便于报告引用）}

\begin{table}[htbp]
  \centering
  \caption[静态 MILP 与仿真对照]{静态 MILP 与 \texttt{fleet\_simulation.py} 主要对照（除上帝视角外）}
  \small
  \begin{tabular}{@{}p{0.28\textwidth}p{0.32\textwidth}p{0.32\textwidth}@{}}
    \toprule
    项目 & 仿真（动态） & 静态 MILP（Gurobi） \\
    \midrule
    任务与充电站几何 & \multicolumn{2}{@{}p{0.64\textwidth}@{}}{可与 \texttt{FleetSimulator} 仅生成任务时对齐（同 \texttt{seed}、同配置）。} \\
    最短路、$v$、$\rho$、$B_{\max}$、$L_{\max}$、$P_{\mathrm{charge}}$ & \multicolumn{2}{@{}p{0.64\textwidth}@{}}{定义一致，用于距离、时间与能耗。} \\
    决策时刻 & 每步 \texttt{step} 在线派车、装批、路径段 & 一次性离线优化（已知全集）。 \\
    车辆路径条数 & 多趟回仓、多批任务 & 当前实现为每车单条闭合路径（单趟抽象）。 \\
    充电排队 & 连续区间 $+$ 槽位预约 $+$ 推进最早开始时刻 & 时间槽上的容量约束（$\le 2$），离散近似。 \\
    得分 & 事件累加 $S$（第~\ref{sec:score}~节） & 代理目标，与 $S$ 不完全同一函数。 \\
    \bottomrule
  \end{tabular}
\end{table}

\subsection{如何撰写对比结论}

建议在报告中分三点写清边界：\textbf{（i）}数据层（任务、仓库、充电站节点、路网与参数）与仿真可对齐；\textbf{（ii）}优化层为静态、单趟抽象 $+$ 离散排队近似，与 \texttt{run()} 的\textbf{多轮动态过程}及\textbf{连续充电预约}不完全同构；\textbf{（iii）}数值对比应以“同实例下的上界/参考解/策略差异说明”为主，避免将 MILP 目标值与仿真 $S$ 直接等同为同一标尺。

\section[真实路网 OSM 与可视化]{真实路网扩展（OpenStreetMap）与“地图”可视化原理}
\label{sec:osm-map}

除网格抽象外，代码中还提供将\textbf{同一套车队模型}（多车、电量、充电排队、动态任务与得分，见前文）迁移到\textbf{OpenStreetMap（OSM）道路几何}上的实现：\texttt{fleet\_osm.py}、\texttt{osm\_graph.py}，以及仅负责拉数导出的 \texttt{osm\_fetch\_demo.py}。本节说明\textbf{使用了哪些库/接口}，以及\textbf{从数据到屏幕}的原理，避免与“嵌入百度/高德瓦片地图”一类方案混淆。

\subsection{使用了什么：无专用地图渲染 SDK}

\begin{itemize}
  \item \textbf{数据来源与协议}：\textbf{OpenStreetMap} 社区数据；通过公共 \textbf{Overpass API} 以 \texttt{POST} 方式查询。请求与解析使用 Python \textbf{标准库} \texttt{urllib}、\texttt{json}，\textbf{不依赖} \texttt{folium}、\texttt{leaflet}、\texttt{osmnx} 等 GIS/前端地图库。
  \item \textbf{图构建与最短路}：在自建的无向加权图上调用与主仿真相同的 \textbf{Dijkstra}（实现于 \texttt{fleet\_simulation.py} 的邻接表版本），同样采用\textbf{按源点缓存距离行}。
  \item \textbf{动态可视化}：使用 Python 标准库 \textbf{\texttt{tkinter}} 的 \texttt{Canvas} 逐帧绘制折线段与图标。\textbf{没有}加载卫星/街道\textbf{瓦片底图}；屏幕上的“地图”是\textbf{矢量路网 $+$ 业务图层}在画布上的几何投影。
  \item \textbf{在线失败回退}：每个 OSM 场景在线请求设置\textbf{总超时 $10\,\mathrm{s}$}。若某一场景首次在线失败（超时/网络/HTTP 错误），当前场景立即回退到 \texttt{osm\_export\_csv/<SCENARIO>\_map\_nodes.csv} 与 \texttt{osm\_export\_csv/<SCENARIO>\_map\_edges.csv}；并且\textbf{同一次运行中后续场景直接走本地 CSV}，避免在 SMALL/MEDIUM/LARGE 上重复等待超时。
\end{itemize}

\subsection[从 Overpass 到路网图]{从 Overpass 响应到图 $G=(V,E)$}

\begin{enumerate}
  \item 在经纬度边界框（bbox：南、西、北、东）内查询带 \texttt{highway=*} 的 \texttt{way}，并使用 \texttt{out geom;} 取得每条道路中心线的\textbf{折线顶点序列}（经纬度）。
  \item 将相邻几何点 $(\lambda_1,\varphi_1)\rightarrow(\lambda_2,\varphi_2)$ 连成\textbf{路段}；对端点坐标做\textbf{固定小数位量化}（如保留到 $10^{-6}$ 度）以合并数值噪声，得到抽象顶点集 $V$。
  \item 每条无向边连接两个量化后的端点，边权取两点间\textbf{大圆距离的球面近似（Haversine）}，记地球半径为 $R$（代码中取 $R\approx 6371000\,\mathrm{m}$），则
  \begin{equation}
    a=\sin^2\frac{\Delta\varphi}{2}+\cos\varphi_1\cos\varphi_2\sin^2\frac{\Delta\lambda}{2},\qquad
    d_{uv}=2R\arcsin(\min\{1,\sqrt{a}\}),
  \end{equation}
  其中 $\Delta\varphi,\Delta\lambda$ 为弧度差。由此\textbf{边权单位为米}，与网格版“抽象距离 $1$”不同；\texttt{SimConfig} 中车速 $v$、单位距离耗电 $\rho$ 在 OSM 模式下解释为 \textbf{m/s} 与 \textbf{每米耗电}，时间轴仍为秒，以保持与 \texttt{step()} 离散推进同一套接口。
  \item 取道路网络的\textbf{最大连通分量}，并在该分量上选取\textbf{仓库}顶点（实现上偏向靠近点集几何中心且度数较高，以减少“死胡同仓”）；任务仅出现在从仓库\textbf{最短路距离有限}的顶点上，保证派发有意义。
\end{enumerate}

\subsection[继承 FleetSimulator]{继承 \texttt{FleetSimulator} 的含义}

\texttt{FleetSimulatorRoad} 以子类方式\textbf{替换图的来源与仓库/候选点初始化}，\textbf{不重写} \texttt{step()}、分段行驶、充电预约与得分累加等核心逻辑。因此 OSM 模式与网格模式在\textbf{控制与事件语义}上一致，差异主要体现在\textbf{拓扑与度量}由规则网格变为\textbf{稀疏平面路图}。

\subsection{动态画面如何把经纬度画到 Tk 画布上}

记当前子图所有顶点经纬度落在经度区间 $[\lambda_{\min},\lambda_{\max}]$、纬度区间 $[\varphi_{\min},\varphi_{\max}]$，画布上预留绘图区宽 $W_p$、高 $H_p$（像素），左下角映射原点像素偏移为 $(o_x,o_y)$。则对任意顶点 $(\lambda,\varphi)$，采用\textbf{线性仿射 $+$ 纵轴翻转}（屏幕 $y$ 向下为正）：
\begin{equation}
  x = o_x + \frac{\lambda-\lambda_{\min}}{\lambda_{\max}-\lambda_{\min}+\varepsilon}\,W_p,\qquad
  y = o_y + \frac{\varphi_{\max}-\varphi}{\varphi_{\max}-\varphi_{\min}+\varepsilon}\,H_p.
\end{equation}
$\varepsilon$ 为防止除零的小常数。路网层：对每条边 $(u,v)\in E$ 连接其两端点在画布上的直线段即可。车辆位置：仿真仍输出最短路顶点序列及每段时间区间；可视化侧在相邻顶点间按\textbf{已行驶弧长占该弧总长}的比例做线性插值（与 \texttt{fleet\_visual.py} 中沿路径插值思想相同，只是把“格子中心坐标”换成上式投影坐标）。

\subsection{小结（与主文网格模型的关系）}

网格版回答“抽象图上的调度原理”；OSM 扩展回答“若把\textbf{真实道路几何}当作图、把\textbf{米}当作距离单位，同一调度模型如何运作以及如何\textbf{自绘}动态演示”。二者\textbf{不共用}同一套 \texttt{rows}$\times$\texttt{cols} 可视化，但\textbf{共用}同一套车队状态机接口。

\section{小结}
\label{sec:summary}

网格图 $+$ Dijkstra（行缓存）、离散时间动态任务、\textbf{仓库多票装车（重量贪心）与最近邻配送序、送完一批再回仓}、分段路径下的充电站槽位与区间预约，以及按时奖励、迟到/过期惩罚、按段路程惩罚构成总分 $S$。上述基线行为以 \texttt{fleet\_simulation.py} 为准。\textbf{元启发式对比}在 \texttt{fleet\_metaheuristic.py} 中实现：装车为同重 EDD 平局；批内序为小规模全排列最优或模拟退火最小化 $J=\alpha_{d}D+\alpha_{\mathrm{late}}L$，与得分系数对齐以便改进访问序时兼顾路程与紧迫度。\textbf{静态 MILP} 在 \texttt{fleet\_gurobi\_static.py} 中实现：与仿真可对齐任务与充电站布局及主要物理参数，但在多趟动态过程、充电排队机制与目标函数上与仿真\textbf{不等价}，详见第~\ref{sec:gurobi-compare}~节。\textbf{真实路网扩展}见第~\ref{sec:osm-map}~节：OSM $+$ Overpass（标准库 HTTP）构图，\texttt{tkinter.Canvas} 经纬度线性投影动态演示（无瓦片地图 SDK）。

\section{电量消耗的具体参数与公式（仅追加说明）}
\label{sec:energy-detail}

本节将代码中的电量逻辑写成显式公式，便于在报告中直接引用。参数与实现对应 \texttt{fleet\_simulation.py}。

\subsection{参数定义（默认场景共用）}
\[
  B_{\max}=500,\quad
  \rho=0.8,\quad
  v=1.0,\quad
  P_{\mathrm{charge}}=75.0.
\]
其中：
\begin{itemize}
  \item $B_{\max}$：电池容量上限（满电值）；
  \item $\rho$：单位距离能耗（电量/距离）；
  \item $v$：行驶速度（距离/时间）；
  \item $P_{\mathrm{charge}}$：充电功率（电量/时间）。
\end{itemize}

\subsection{单段行驶：能耗与时间}
对任意一段最短路距离 $d$，行驶时间与能耗分别为
\begin{equation}
  t_{\mathrm{travel}}=\frac{d}{v},\qquad
  E(d)=\rho d.
\end{equation}
若段起点电量为 $B_{\mathrm{start}}$，且可直达，则段终点电量
\begin{equation}
  B_{\mathrm{end}}=B_{\mathrm{start}}-\rho d,\qquad
  \rho d\le B_{\mathrm{start}}.
\end{equation}

\subsection{直达判定与“先充后送”判定}
设当前段为 $u\rightarrow w$，其最短路距离为 $d_{uw}$。
\begin{itemize}
  \item \textbf{直达条件}：$\rho d_{uw}\le B_{\mathrm{now}}$。
  \item \textbf{若不满足直达}：尝试经某充电站 $c$，并要求
  \begin{equation}
    \rho d_{uc}\le B_{\mathrm{now}},\qquad
    \rho d_{cw}\le B_{\max}.
  \end{equation}
  第一式保证“当前电量可到站”，第二式保证“满电后可到目标”。
\end{itemize}

\subsection{充电时长与充后电量}
车辆到站时电量
\begin{equation}
  B_{\mathrm{arr}}=B_{\mathrm{now}}-\rho d_{uc}.
\end{equation}
充至满电所需能量缺口与时长
\begin{equation}
  \Delta B=\max\{0,\ B_{\max}-B_{\mathrm{arr}}\},\qquad
  t_{\mathrm{charge}}=\frac{\Delta B}{P_{\mathrm{charge}}}.
\end{equation}
充电结束后按实现取
\begin{equation}
  B_{\mathrm{after}}=B_{\max}.
\end{equation}

\subsection{排队与开始充电时刻}
设到站时刻为 $t_{\mathrm{arrive}}$。充电站并发槽位数为 $2$，开始充电时刻为
\begin{equation}
  t_{\mathrm{start}}
  =\min\{t\ge t_{\mathrm{arrive}}:\ \text{该时刻占用会话数}<2\}.
\end{equation}
因此结束时刻
\begin{equation}
  t_{\mathrm{end}}=t_{\mathrm{start}}+t_{\mathrm{charge}}.
\end{equation}

\subsection{经站补能后的到达时刻与终点电量}
若选择 $u\rightarrow c\rightarrow w$，则到达目标时刻
\begin{equation}
  t_{\mathrm{arrive},w}
  =t_0+\frac{d_{uc}}{v}+\left(t_{\mathrm{start}}-t_{\mathrm{arrive}}\right)+t_{\mathrm{charge}}+\frac{d_{cw}}{v}.
\end{equation}
目标点电量
\begin{equation}
  B_w=B_{\max}-\rho d_{cw}.
\end{equation}

\subsection{充电站选择准则（代码实现口径）}
当直达不可行时，在所有可行充电站集合 $\mathcal{C}_{\mathrm{feasible}}$ 中按\textbf{字典序}选站：
\begin{equation}
  c^*=\arg\min_{c\in\mathcal{C}_{\mathrm{feasible}}}
  \Bigl(
    \underbrace{\mathbf{1}[\text{arrive\_busy}(c)\ge \text{slots}(c)]}_{\text{是否无空位（有空位优先）}},
    \underbrace{d_{uc}}_{\text{到站距离}},
    \underbrace{t_{\mathrm{start}}(c)}_{\text{最早可开始充电时刻}},
    \underbrace{\text{id}(c)}_{\text{稳定打破平局}},
    \underbrace{t_{\mathrm{arrive\_target}}(c)}_{\text{可行性兜底比较}}
  \Bigr).
\end{equation}
其中第一优先级体现“\textbf{有无位置}”，第二优先级体现“\textbf{充电点距离}”。在前两级相同情况下，再比较可开始充电时刻与到达目标时刻。到达时刻仍按
\begin{equation}
  t_{\mathrm{arrive\_target}}(c)=
  t_0+\frac{d_{uc}}{v}
  +\text{queue\_wait}(c,t_{\mathrm{arrive}})
  +\frac{\max(0,\ B_{\max}-B_{\mathrm{arr}})}{P_{\mathrm{charge}}}
  +\frac{d_{cw}}{v}.
\end{equation}
因此实现不再是“单纯最近站”或“仅最早到达”单指标，而是与你的实验口径一致的\textbf{空位优先 + 距离优先}。

\end{document}
